\begin{table}[htbp]
\centering
\caption{Descomposición del cambio de la razón SHRFSPF/PIB en cuatro variables (puntos del PIB)}
\label{tab:C12_4}
\small
\begin{tabular}{lSSSSSSSSSSS}
\toprule
{Caso} & {$d_{t-1}$} & {$d_t$} & {Obs.} & {RFSPF} & {Crec.} & {Infl. denom.} & {Compens.} & {Infl. neta} & {Tipo cambio} & {Suma} & {Residuo} \\
\midrule
Prueba 2022, añada del CGPE 2024 (la validada en la rubrica 2023) & \num{49.2} & \num{47.7} & \num{-1.5} & \num{4.3} & \num{-1.73} & \num{-3.09} & \num{2.97} & \num{-0.12} & \num{-0.8} & \num{-1.32} & \num{-0.18} \\
Prueba 2022, añada que nombra la instruccion 2025 (50.7 -> 49.4) & \num{50.7} & \num{49.4} & \num{-1.3} & \num{4.5} & \num{-1.78} & \num{-3.18} & \num{3.06} & \num{-0.12} & \num{-0.82} & \num{-1.29} & \num{-0.01} \\
Aplicacion a 2025 (SHRFSPF 51.4 -> 51.4) & \num{51.4} & \num{51.4} & \num{0.0} & \num{3.9} & \num{-1.06} & \num{-2.12} & \num{2.07} & \num{-0.05} & \num{-0.79} & \num{-0.07} & \num{0.07} \\
\bottomrule
\end{tabular}
\par\vspace{2pt}\footnotesize Fuente: construcción propia del ITED con el CGPE 2025. Los dos primeros renglones son la prueba de aceptación del marco sobre 2022, con las dos añadas que estaban en disputa.
\end{table}
